% !TEX root = ../../tir_monograph_v12.tex
% V12_SOURCE_MAP:
% - TIR/monograph/chapters/ch03_l_constants.tex
% - TIR/monograph/chapters/ch04_quark_primes.tex
% - TIR/STRUCTURAL_CHOICES.md
% - TIR/foundations/TIR_PLATONIC_L_CONSTANTS_CLOSURE_V0_1.md
% Validators:
% - TIR/validation/tir_v12_discrete_labels_audit_v0_1.py
% - TIR/validation/tir_platonic_l_constants_closure_v0_1.py
\chapter[Discrete Structural Labels]{Discrete Structural Labels -- $L_i$, Primes and Provenance}
\label{ch:v12_discrete_labels}

\paragraph{Established context.}
The Collatz map, prime-number sequence, Platonic classification and finite groups
\(A_4,S_4,A_5\) are standard mathematical objects.  TIR uses selected finite
properties of those objects inside an explicitly typed structural layer.  The
v12.1 synchronization separates three levels: standard mathematical facts, the
TIR selection/closure rule, and any downstream physical interpretation.

\section{Platonic index alphabet}

For a convex regular polyhedron with Schlaefli symbol \(\{p,q\}\),
\[
\frac1p+\frac1q>\frac12,
\qquad p,q\ge3.
\]
The integer solutions are exactly
\[
\boxed{
\{3,3\},\{3,4\},\{3,5\},\{4,3\},\{5,3\}.
}
\]
Thus the complete index alphabet appearing in the convex Platonic
classification is
\[
\boxed{\{3,4,5\}}.
\]
The triangular-face branch is
\[
\{3,3\},\qquad\{3,4\},\qquad\{3,5\},
\]
with orientation-preserving rotational groups
\[
G_3\simeq A_4,
\qquad
G_4\simeq S_4,
\qquad
G_5\simeq A_5,
\]
and orders
\[
|A_4|=12,
\qquad |S_4|=24,
\qquad |A_5|=60.
\]
Equivalently, the spherical triangle groups \((2,3,q)\) have
\[
|G_q|
=\frac{2}{\frac12+\frac13+\frac1q-1},
\qquad q\in\{3,4,5\}.
\]
At \(q=6\), the spherical excess vanishes; hence the finite spherical branch
terminates at \(q=5\).

\section{Common tetrahedral root and subgroup indices}

The tetrahedral rotational group satisfies
\[
A_4\triangleleft S_4,
\]
so
\[
\boxed{[S_4:A_4]=\frac{24}{12}=2}.
\]
A point stabilizer in the natural five-point action of \(A_5\) is isomorphic to
\(A_4\), giving
\[
\boxed{[A_5:A_4]=\frac{60}{12}=5}.
\]
Define the two extension coset spaces over the tetrahedral root,
\[
X_4:=S_4/A_4,
\qquad
X_5:=A_5/A_4.
\]
Then
\[
|X_4|=2,
\qquad
|X_5|=5.
\]
TIR identifies
\[
\boxed{L_4:=|X_4|=2},
\qquad
\boxed{L_5:=|X_5|=5}.
\]

\section{TIR root-extension closure and $L_3$}

The TIR structural step is explicit:

\begin{quote}
For the tetrahedral root, the complete non-self-dual Platonic extension carrier
is the disjoint union of all admissible extension coset spaces.
\end{quote}

Therefore
\[
X_3^{\rm closure}:=X_4\sqcup X_5
=S_4/A_4\sqcup A_5/A_4.
\]
Since the union is disjoint,
\[
|X_3^{\rm closure}|=|X_4|+|X_5|=2+5=7,
\]
and TIR defines
\[
\boxed{L_3:=|X_3^{\rm closure}|=7}.
\]
Thus
\[
\boxed{
(L_3,L_4,L_5)
=\left(
|S_4/A_4|+|A_5/A_4|,
|S_4/A_4|,
|A_5/A_4|
\right)
=(7,2,5).
}
\]
\vTwelveStatus{B}{--}{PASS}

The standard finite-group facts are not fitted.  The root-extension closure
rule is a TIR structural definition; the triple follows exactly once that rule
is admitted.  The finite-group derivation alone does not establish a physical
law, a dynamical role for the labels, or a derivation from a specific Kaehler
manifold.

Canonical theorem:
\path{TIR/foundations/TIR_PLATONIC_L_CONSTANTS_CLOSURE_V0_1.md}.

Deterministic validator:
\path{TIR/validation/tir_platonic_l_constants_closure_v0_1.py}.

\section[Independent finite Collatz crosscheck for L3]{Independent finite Collatz crosscheck for $L_3$}

Define the usual Collatz step on positive integers by
\begin{equation}
C(n)=
\begin{cases}
3n+1, & n\ \text{odd},\\
n/2, & n\ \text{even}.
\end{cases}
\label{eq:v12-collatz-map}
\end{equation}
For the single seed required here, direct iteration gives the finite orbit
\begin{equation}
\boxed{
3\to10\to5\to16\to8\to4\to2\to1.
}
\label{eq:v12-collatz-three-orbit}
\end{equation}
There are exactly seven steps before the first occurrence of \(1\).  Hence the
pre-existing arithmetic provenance gives independently
\begin{equation}
\boxed{\operatorname{depth}_{\rm Collatz}(3)=7=L_3.}
\label{eq:v12-L3}
\end{equation}
\vTwelveStatus{B}{--}{PASS}
The validator evaluates only this displayed finite orbit; no global assertion
about all positive Collatz seeds enters the calculation.

\section[Independent twin-prime crosschecks for L4 and L5]{Independent twin-prime crosschecks for $L_4$ and $L_5$}

From the prime pair \((3,5)\), the older arithmetic route gives
\begin{equation}
\boxed{5-3=2=L_4,}
\qquad
\boxed{5=L_5.}
\label{eq:v12-L4-L5}
\end{equation}
\vTwelveStatus{B}{--}{PASS}
These identities are retained as independent crosschecks of the values obtained
from the Platonic subgroup-index construction, not as inputs used to fit that
construction.

Useful downstream combinations include
\[
L_3+L_4=9,
\qquad
L_3-1=6,
\qquad
\frac{L_4}{L_5}=\frac25,
\qquad
L_4^{L_5}=2^5=32.
\]
Additional exact checks connecting the Platonic index alphabet and the three
rotation-group orders are
\[
3^2+4^2=5^2,
\qquad
3+4+5=12=|A_4|,
\]
\[
2(3+4+5)=24=|S_4|,
\qquad
3\cdot4\cdot5=60=|A_5|.
\]
These are consistency checks, not extra derivation premises.

\section{Prime candidate set for quark labels}

The stronger derivation of the \(L_i\) values does not automatically derive the
quark-prime map.  That map retains its own provenance and identifiability gate.
Reserve the first prime \(p_1=2\) for the binary structural branch and define the
ordered candidate list
\begin{equation}
\mathcal P_Q
=(p_2,p_3,p_4,p_5,p_6,p_7)
=(3,5,7,11,13,17).
\label{eq:v12-prime-candidate-list}
\end{equation}
Let the ordered quark-slot list be
\begin{equation}
\mathcal F_Q=(u,d,s,c,b,t).
\label{eq:v12-quark-slot-list}
\end{equation}

\section{Arithmetic constraints and residual ambiguity}

The historical assignment satisfies
\begin{equation}
(u,d,s,c,b,t)=(3,5,7,11,13,17),
\label{eq:v12-quark-prime-assignment}
\end{equation}
together with
\[
d-u=2,
\qquad
s=L_3=7,
\qquad
|s-c|=4,
\qquad
|b-t|=4.
\]
An exhaustive scan over all
\[
6!=720
\]
bijections between \(\mathcal F_Q\) and \(\mathcal P_Q\) shows that these
arithmetic conditions alone leave exactly two survivors:
\begin{align}
(u,d,s,c,b,t)&=(3,5,7,11,13,17),\\
(u,d,s,c,b,t)&=(3,5,7,11,17,13).
\end{align}
Thus the displayed arithmetic relations leave one residual heavy-pair exchange
\(b\leftrightarrow t\).
\vTwelveStatus{D}{--}{PASS}

\section{Typed ordering rule}

The v12 structural selection is the explicit order-preserving map
\begin{equation}
\boxed{
q:\mathcal F_Q\to\mathcal P_Q,
\qquad
q((\mathcal F_Q)_i)=(\mathcal P_Q)_i,
\quad i=1,\ldots,6.
}
\label{eq:v12-quark-prime-typed-rule}
\end{equation}
Equivalently,
\[
q(u)=3,\quad q(d)=5,\quad q(s)=7,\quad q(c)=11,\quad q(b)=13,\quad q(t)=17.
\]
Within the declared ordered candidate set and ordered flavour slots, this rule
has exactly one assignment.  The uniqueness is uniqueness of the typed
selection rule; a scheme- and scale-defined physical quark-mass map remains a
separate open gate.
\vTwelveStatus{B}{--}{PASS}

\section{Reproducibility receipt}

The deterministic audit
\path{TIR/validation/tir_v12_discrete_labels_audit_v0_1.py}
checks the finite Collatz orbit and enumerates all \(720\) prime assignments.
The independent validator
\path{TIR/validation/tir_platonic_l_constants_closure_v0_1.py}
enumerates the five Platonic Schlaefli pairs, constructs the finite permutation
groups and cosets, verifies the subgroup indices and reproduces
\((L_3,L_4,L_5)=(7,2,5)\) without using the older arithmetic values as inputs.
